\chapter{Seeley--DeWitt Coefficients at Working Thickness}
\label{app:seeley-dewitt}

This appendix records the heat-kernel formulae used in Chapter~8.
It introduces no new physical hypothesis. Conventions follow
Section~\ref{sec:ch2-heatkernel}: Euclidean section first,
Lorentzian continuation afterwards, signature of the main text
restored only when an Einstein--Hilbert term is isolated.

\section{Laplace-type operators}
\label{sec:appA-laplace}

Let \(V\to M\) be a vector bundle over a smooth four-manifold with
Riemannian metric \(g_{\mu\nu}\). A second-order differential
operator on sections of \(V\) is of Laplace type if it can be
written
\begin{equation}
  \label{eq:appA-O}
  \mathcal{O}
  =-g^{\mu\nu}D_\mu D_\nu+E,
\end{equation}
where \(D_\mu=\nabla_\mu+\mathcal{A}_\mu\) is a connection on \(V\)
(including the Levi-Civita connection on tensor factors) and \(E\)
is a bundle endomorphism. The curvature of \(D\) is
\begin{equation}
  \label{eq:appA-F}
  \mathcal{F}_{\mu\nu}
  =[D_\mu,D_\nu].
\end{equation}
On the locked condensate block of Chapter~5, \(V\) is the bundle of
tensors \(\delta\Phi_{\mu\nu}\) together with whatever ghosts are
included in a given computation.

The heat operator \(e^{-s\mathcal{O}}\) is defined for \(s>0\) by
the usual functional calculus after \(\mathcal{O}\) has been made
positive by a Wick rotation. Its integral kernel
\(K(s;x,y)\) satisfies
\begin{equation}
  \label{eq:appA-heat}
  \partial_s K+\mathcal{O}_x K=0,
  \qquad
  K(0;x,y)=\delta(x,y).
\end{equation}

\section{Small-time expansion}
\label{sec:appA-expand}

On the diagonal in four dimensions,
\begin{equation}
  \label{eq:appA-SDW}
  \mathrm{tr}_V K(s;x,x)
  \sim
  \frac{1}{(4\pi s)^{2}}
  \sum_{k=0}^{\infty}
  s^k a_k(x),
  \qquad
  s\to 0^+.
\end{equation}
Each \(a_k\) is a local scalar density of mass dimension \(2k\),
built from \(R_{\mu\nu\rho\sigma}\), \(\mathcal{F}_{\mu\nu}\),
\(E\), and their covariant derivatives. The first three
coefficients, up to total derivatives that do not enter the bulk
effective action, are
\begin{equation}
  \label{eq:appA-a0}
  a_0=\mathrm{tr}_V\mathbb{I},
\end{equation}
\begin{equation}
  \label{eq:appA-a1}
  a_1
  =\mathrm{tr}_V
   \Bigl(
     \tfrac16 R-E
   \Bigr),
\end{equation}
\begin{equation}
  \label{eq:appA-a2}
  \begin{aligned}
    a_2
    =\mathrm{tr}_V
     \Biggl(
       &\frac{1}{180}
        \bigl(
          R_{\mu\nu\rho\sigma}R^{\mu\nu\rho\sigma}
          -R_{\mu\nu}R^{\mu\nu}
          +\tfrac52 R^2
        \bigr)
        +\frac12 E^2
        -\frac16 RE\\
       &+\frac{1}{12}\mathcal{F}_{\mu\nu}\mathcal{F}^{\mu\nu}
        +\frac{1}{30}D^2 R
        -\frac16 D^2 E
     \Biggr).
  \end{aligned}
\end{equation}
Different references redistribute the \(D^2 R\) and \(D^2 E\) terms
by integration by parts. For the bulk action of Chapter~8 those
terms are discarded.

The numerical coefficients in~\eqref{eq:appA-a2} are the standard
DeWitt values for a Laplace-type operator. They do not determine
the coefficients \(A,B,C\) in the conditional
parameterization~\eqref{eq:Geff-structure}. Their evaluation requires
the specified operators, fibre traces, determinant weights, and
thresholds; this remains part of Problem~F in Chapter~17.

\section{Proper-time integrals}
\label{sec:appA-integrals}

For the sectors actually included in a gauge-fixed calculation, let
\(\mathcal{S}_{\mathrm{sec}}\) be their index set, \(w_s\) their determinant weights
relative to a real bosonic determinant, and \(\mathcal O_s\) their
operators. Define
\begin{equation}
  \label{eq:appA-sector-sum}
  W_E^{(1)}
  =\frac12\sum_{s\in\mathcal{S}_{\mathrm{sec}}}w_s
        \operatorname{Tr}_{V_s}\ln\mathcal O_s,
  \qquad
  \mathcal N_0
  :=\sum_{s\in\mathcal{S}_{\mathrm{sec}}}w_s
        \operatorname{tr}_{V_s}\mathbb I_s .
\end{equation}
The corresponding proper-time representation is
\begin{equation}
  \label{eq:appA-W}
  W_E^{(1)}
  =-\frac12\sum_{s\in\mathcal{S}_{\mathrm{sec}}}w_s
    \int_{s_{\min,s}}^{s_{\max,s}}\frac{\mathrm{d}s}{s}\,
    \operatorname{Tr}_{V_s}e^{-s\mathcal O_s}.
\end{equation}
The expressions below are term-by-term integrals of the truncated,
unresummed small-\(s\) heat-kernel expansion through \(a_2\), not
exact finite-cutoff integrations of massive determinants. If \(E_s\)
contains \(m_s^2\), resumming \(\exp(-s m_s^2)\) replaces polynomial
cutoff terms by sector-dependent threshold functions. A common pair of
limits applies only to sectors explicitly assigned those limits;
otherwise the sector-dependent limits are retained. Inserting the
heat-kernel expansion and integrating term by term gives
\begin{equation}
  \label{eq:appA-Wexp}
  \begin{aligned}
      W_E^{(1)}
      &\supset
      \sum_{s\in\mathcal{S}_{\mathrm{sec}}}w_s
      \int\mathrm{d}^4x\,\sqrt{g}\,
      \Biggl[
        -\frac{a_{0,s}}{64\pi^2}
          ({\LUV}_s^4-\mu_s^4)
        -\frac{a_{1,s}}{32\pi^2}
          ({\LUV}_s^2-\mu_s^2)\\
      &\qquad\qquad
        -\frac{a_{2,s}}{32\pi^2}
          \ln\Bigl(\frac{{\LUV}_s^2}{\mu_s^2}\Bigr)
      \Biggr],
  \end{aligned}
\end{equation}
where \(a_{k,s}\) are the traced coefficients on \(V_s\). With common
limits explicitly assigned to every included sector, the \(a_0\) term
is \(-\mathcal N_0(\LUV^4-\mu^4)/(64\pi^2)\).

For the Lorentzian vacuum-sign statement adopt \(S_E=-iS_L\) under
\(t=-i\tau\) and define the Lorentzian vacuum term by
\(S_L\supset-\rho_\Lambda\int\mathrm{d}^4x\,\sqrt{-g}\).
For one real scalar, \(a_0=1\), and only under these conventions,
\(W_E/\mathrm{Vol}\supset-(\LUV^4-\mu^4)/(64\pi^2)\) corresponds to
\(\rho_\Lambda=-(\LUV^4-\mu^4)/(64\pi^2)\). The common-limit sector
sum gives \(\rho_\Lambda=-\mathcal N_0(\LUV^4-\mu^4)/(64\pi^2)\)
only when all displayed sectors share the assigned limits, within the
same truncated, unresummed expansion. The Lorentzian sign is
conditional on this continuation/action convention.

\section{Worked example: a non-minimal scalar}
\label{sec:appA-scalar}

For a single real scalar with
\begin{equation}
  \label{eq:appA-scalar}
  \mathcal{O}=-\nabla^2+\xi R+m^2,
\end{equation}
one has \(E=\xi R+m^2\), \(\mathcal{F}_{\mu\nu}=0\), and
\(\mathrm{tr}_V\mathbb{I}=1\). Then
\begin{equation}
  \label{eq:appA-scalar-a1}
  a_1=\bigl(\tfrac16-\xi\bigr)R-m^2.
\end{equation}
The curvature piece of \(a_1\) cancels at \(\xi=1/6\). That
cancellation removes this scalar's contribution to the curvature
coefficient through \(a_1\). The NAMCG radial
mode is not assumed to sit at \(\xi=1/6\); Exercise~3 of
Chapter~8 is the corresponding warning.

For one real scalar the Euclidean \(a_0\) contribution is
\(-(\LUV^4-\mu^4)/(64\pi^2)\). Under the continuation and Lorentzian
vacuum-action convention stated above, this corresponds to
\(\rho_\Lambda=-(\LUV^4-\mu^4)/(64\pi^2)\). For a collection of
sectors the total density includes their determinant weights and
sector-dependent thresholds.

\section{Fibre traces on the locked block}
\label{sec:appA-fibre}

If \(V\) is the space of symmetric traceless tensors on a
four-manifold, \(\mathrm{tr}_V\mathbb{I}=9\). If the algebraic
split of Chapter~5 is retained before gauge fixing, the sixteen
components of \(\delta\Phi_{\mu\nu}\) give
\(\mathrm{tr}_V\mathbb{I}=16\). Ghosts for diffeomorphisms
contribute with the opposite sign and with their own Laplace-type
operators. Any local-Lorentz ghost sector must be specified together
with the gauge fixing of the independent Lorentz connection;
Theorem~\ref{thm:locking} does not supply that operator.

\(\mathcal N_0\) is the weighted fibre trace over the sectors actually
included, not automatically the locked-block dimension or a
physical-mode count. A total locked-block result must specify the
gauge-fixed fluctuation, Goldstone, diffeomorphism-ghost and
local-Lorentz-ghost operators, determinant weights, and assigned
thresholds. No factor of \(1/6\) is absorbed into \(\mathcal N_0\);
the complete calculation is Problem~F.

\section{What this appendix will not do}
\label{sec:appA-not}

It will not list \(a_3\) and \(a_4\). Those coefficients generate
six- and eight-derivative operators that the book treats as
suppressed by further powers of \(\LUV^{-1}\).

It will not prove the coefficients~\eqref{eq:appA-a0}--\eqref{eq:appA-a2}.
Standard derivations (coincidence limits of the Synge world
function, or the recursive DeWitt transport equations) are in the
references collected in \texttt{bib/namcg.bib}. The book uses the
coefficients; it does not redo the recursion.

It will not evaluate \(a_2\) on the Einsteinian cubic background
of Chapter~13. That evaluation is part of Problem~J if it is
needed, not part of the working thickness of Chapter~8.

\section{Exercises}
\label{sec:appA-exercises}

\begin{enumerate}
  \item \textbf{Dimensions.}
    Assign mass dimensions to \(s\), \(a_0\), \(a_1\), \(a_2\),
    and \(\mathrm{tr}_V K(s;x,x)\). Confirm that each term in
    \eqref{eq:appA-Wexp} has dimension four.
  \item \textbf{Conformal scalar.}
    Set \(\xi=1/6\) and \(m=0\) in~\eqref{eq:appA-scalar-a1}.
    What remains of \(a_1\)? What remains of the \(\LUV^2 R\)
    term in \(W_E\)?
  \item \textbf{Gauss--Bonnet combination.}
    Rewrite the curvature part of~\eqref{eq:appA-a2} in the
    basis \(\{C_{\mu\nu\rho\sigma}C^{\mu\nu\rho\sigma},R^2,
    E_{\mathrm{GB}}\}\). Which combination multiplies a
    topological density in four dimensions?
  \item \textbf{Ghost sign.}
    A diffeomorphism ghost in the background-field gauge is a
    complex vector with a Laplace-type operator. What factor
    multiplies its contribution to \(a_0\) relative to a real
    scalar, and why is the sign opposite to a physical field?
  \item \textbf{Conditional curvature coefficient.}
    Given specified sector operators, determinant weights, proper-time
    limits, and a finite-counterterm scheme, derive the truncated
    Euclidean curvature coefficient from \(a_1\) and \(a_2\). Identify
    the additional data needed for its Lorentzian continuation and
    comparison with \(G_{\mathrm{eff}}\).
\end{enumerate}